\chapter*{Quick Revision: Frequently Asked Facts}
\addcontentsline{toc}{chapter}{Quick Revision: Frequently Asked Facts}
\markboth{Quick Revision}{Quick Revision}

\small
\begin{multicols}{2}
\subsection*{Teaching}
\begin{itemize}
\item Levels of teaching: \textbf{Memory} (Herbart) $\rightarrow$ \textbf{Understanding} (Morrison) $\rightarrow$ \textbf{Reflective} (Hunt).
\item Evaluation: \textbf{diagnostic} (find difficulties), \textbf{formative} (during teaching, improve learning), \textbf{summative} (at the end, certify).
\item Bloom's cognitive levels: remember, understand, apply, analyse, evaluate, create.
\item Maxims: known to unknown, simple to complex, concrete to abstract, particular to general (induction).
\item Presage (before) -- Process (during) -- Product (after) model of teaching.
\end{itemize}
\subsection*{Research}
\begin{itemize}
\item Steps: problem $\rightarrow$ review $\rightarrow$ hypothesis $\rightarrow$ design and sample $\rightarrow$ tool $\rightarrow$ data $\rightarrow$ analysis $\rightarrow$ report.
\item Probability sampling: simple random, systematic, stratified, cluster, multistage. Non-probability: purposive, quota, convenience, snowball.
\item Normal curve: $\pm1\sigma=68.26\%$, $\pm2\sigma=95.44\%$, $\pm3\sigma=99.73\%$.
\item Scales: nominal, ordinal, interval, ratio.
\item Two group means: $t$-test; three or more: ANOVA ($F$); frequencies: $\chi^2$.
\item Historical research: external criticism (authenticity), internal criticism (credibility).
\end{itemize}
\subsection*{Communication}
\begin{itemize}
\item Kinesics (body), Proxemics (space), Haptics (touch), Chronemics (time), Paralanguage (voice), Oculesics (eyes).
\item Sender encodes, receiver decodes; feedback completes the cycle.
\item Barriers: semantic, psychological, physical (noise), organisational, cultural.
\item 7 Cs: clarity, conciseness, concreteness, correctness, coherence, completeness, courtesy.
\end{itemize}
\subsection*{Indian logic}
\begin{itemize}
\item Pramanas: pratyaksha, anumana, upamana, shabda, arthapatti, anupalabdhi.
\item Panchavayava: pratijna, hetu, udaharana, upanaya, nigamana.
\item Vyapti: invariable concomitance of hetu (smoke) and sadhya (fire).
\end{itemize}
\columnbreak
\subsection*{ICT}
\begin{itemize}
\item 1 byte $=$ 8 bits; KB $<$ MB $<$ GB $<$ TB (factor 1024).
\item URL, HTTP, HTML, FTP, IP address, LAN $<$ MAN $<$ WAN.
\item SWAYAM (MOOCs, 2017), SWAYAM Prabha (40 DTH channels), NPTEL, e-PG Pathshala, Shodhganga, NDLI, NAD, DIKSHA.
\end{itemize}
\subsection*{Environment}
\begin{itemize}
\item AQI bands: 0--50 good, 51--100 satisfactory, 101--200 moderate, 201--300 poor, 301--400 very poor, 401--500 severe.
\item Stockholm 1972, Montreal Protocol 1987, Rio Earth Summit / Agenda 21 1992, Kyoto 1997 (Doha Amendment 2012), Rio+20 2012, Paris 2015.
\item 17 SDGs: 4 Quality education, 5 Gender equality, 6 Water, 7 Energy, 13 Climate, 14 Life below water, 15 Life on land.
\item World Environment Day: 5 June; Ozone Day: 16 September.
\end{itemize}
\subsection*{Higher education}
\begin{itemize}
\item UGC 1956; AICTE 1987 (statutory); NCTE 1993; NAAC 1994 (accreditation valid 5 years); NIRF 2015/16.
\item Radhakrishnan Commission 1948--49; Kothari Commission 1964--66 (10+2+3, 6\% of GDP); NPE 1986, POA 1992; NEP 2020 (5+3+3+4, ABC, HECI with NHERC, NAC, HEGC, GEC).
\item National Education Day: 11 November (Maulana Azad); Teachers' Day: 5 September.
\item CBCS: 1 credit $=$ 1 hour of lecture per week per semester.
\end{itemize}
\end{multicols}
\normalsize

\vfill
\begin{center}
\begin{tcolorbox}[colback=soft,colframe=brand,width=0.85\linewidth,halign=center]
\sffamily\large All the best for your GSET examination!\\[4pt]
\normalsize\rmfamily Found an error or have a suggestion? Please share it with the author\\ so that future editions can be improved.
\end{tcolorbox}
\end{center}
